\paragraph{Spatial consistency of the same coupled action.}
The dressed scalar construction admits an action and first-variation
estimate on the original solid under spatial refinement. The estimate uses
the same edge integrals, Whitney potentials and potential-dependent scalar
basis, rather than replacing them by a different finite-element family.
Whitney interpolation and its commuting exterior derivative are standard
\cite{ArnoldFalkWinther2010}; the additional point here is to control the
nonlinear dressing and its full derivative in this particular action.

Let \(\mathcal T_0\) be the twenty-tetrahedron cone, \(I=[0,T]\),
and let \(\mathcal T_\delta\) be conforming refinements respecting every
macro-face, with maximum diameter \(\delta\leq1\) and a common bound
on diameter divided by inradius. Time remains continuous. Give
\(U=(A,\phi,\psi)\) the norm
\[
 \|U\|_X=\max_{K\in\mathcal T_0}
 \bigl(\|U\|_{L^\infty(I;W^{2,\infty}(K))}
       +\|\partial_tU\|_{L^\infty(I;W^{2,\infty}(K))}\bigr).
\]
Use continuous representatives on each closed macro-cell and require
\(\phi,\psi\in H^1(\Omega)\) and \(A\in H(\operatorname{curl};\Omega)\)
at each time, with the same trace conditions on their time derivatives.
The componentwise interpretation applies to complex \(\psi\).
Let \(a_{ij}=\int_{v_i}^{v_j} A\cdot dx\), take nodal samples of
\(\phi,\psi\), and write \(A_\delta=I_1A\),
\(\phi_\delta=I_0\phi\),
\(\Psi_\delta=\sum_i\lambda_i e^{ie\theta_i}\psi(v_i)\), where
\(\theta_i=\int_{v_i}^{x} I_1A\cdot dx=\sum_j a_{ij}\lambda_j\).
These paths integrate the interpolated potential, not the original smooth
potential. Continuous tangential traces agree on shared faces, hence on
their edges; connectivity of the cells incident to an edge makes its
oriented integral single valued. Scalar face traces agree as well.
No global continuity of the normal component of \(A\) is required.
Define \(S_\delta(U)\) by exact space and time integration
of \eqref{eq:whitney-charged-action} with these fields. Define \(S(U)\)
by the same continuum scalar-electrodynamics density with \(A,\phi,\psi\),
\(E=-\partial_tA-\nabla\phi\) and \(B=\operatorname{curl}A\).

\begin{theorem}[Smooth-window action consistency]
\label{thm:whitney-spatial-consistency}
Fix finite \(T,R,e,m^2,g\), the macro-mesh and a shape-regularity bound.
There is a constant \(C_R\), independent of the refinement, such that
for every \(\|U\|_X\leq R\),
\begin{equation}
 |S_\delta(U)-S(U)|+
 \sup_{\substack{V\in X\\\|V\|_X\leq1}}
 |DS_\delta(U)[V]-DS(U)[V]|\leq C_R\delta.
 \label{eq:whitney-spatial-action-bound}
\end{equation}
Variations obey the same trace requirements. The derivative is the full
real derivative with respect to all three sampled fields, including the
dependence of the scalar basis on the edge potential. This estimate is a
consistency statement on bounded smooth test fields; it does not assert
convergence of nonlinear trajectories or of arbitrary finite-energy data.
\end{theorem}

\begin{proof}
All following bounds are uniform on each small tetrahedron of diameter
\(h\), with constants depending only on the displayed data. Shape regularity
gives \(|\nabla\lambda_i|\leq C/h\). The usual elementary Taylor
estimates for nodal interpolation give
\(\|I_0f-f\|_\infty\leq Ch^2\|f\|_{W^{2,\infty}}\) and
\(\|\nabla(I_0f-f)\|_\infty\leq Ch\|f\|_{W^{2,\infty}}\).
Whitney edge interpolation reproduces constant vector fields, is locally
bounded in the sup norm, and commutes with the exterior derivative by
Stokes's theorem. Consequently
\[
 \|I_1A-A\|_\infty+
 \|\operatorname{curl}(I_1A-A)\|_\infty
 \leq Ch\|A\|_{W^{2,\infty}}.
\]
For clarity, the curl estimate also follows without a general projection
theorem: subtract the affine Taylor polynomial of \(A\). Both its curl
and the interpolated curl are constant; their fluxes through every face
agree by Stokes's theorem and exact sampled edge circulations. The face
normals span \(\mathbb R^3\), so the curls agree. The remaining edge
integrals are \(O(h^3)\), while curls of edge basis functions are
\(O(h^{-2})\). The same estimates apply to \(\partial_tA\), and
to a variation \(V=(B,\eta,\zeta)\), by linearity. Thus the electric
and magnetic fields and their variations have \(O(h)\) errors.

For the scalar field, antisymmetry of \(a_{ij}\) gives the exact
cancellation
\[
 \sum_i\lambda_i\theta_i=\sum_{ij}\lambda_i a_{ij}\lambda_j=0.
\]
The same identity holds after differentiating in space, time or a potential
direction. Put \(r_i=\psi(v_i)-\psi(x)\) and
\(q(s)=e^{ies}-1-ies\). The exact error decomposition is
\[
 \Psi_\delta-I_0\psi
 =\psi(x)\sum_i\lambda_i q(\theta_i)
   +\sum_i\lambda_i(e^{ie\theta_i}-1)r_i.
\]
Here \(\theta_i,\partial_t\theta_i,r_i,\partial_tr_i=O(h)\),
their spatial gradients are bounded, and
\(q(\theta_i)=O(h^2)\), \(q'(\theta_i)=O(h)\).
Differentiating the displayed identity gives
\[
 \|\Psi_\delta-\psi\|_\infty+
 \|\partial_t(\Psi_\delta-\psi)\|_\infty\leq C_Rh^2,
 \qquad
 \|\nabla(\Psi_\delta-\psi)\|_\infty\leq C_Rh.
\]
In particular the time estimate retains the phase velocity, and does not
hold by simply deleting it from the definition.

To check that differentiation in the fields preserves these estimates,
write \(b_{ij}=\int_{v_i}^{v_j}B\cdot dx\) and
\(\beta_i=\sum_j b_{ij}\lambda_j\). The full derivative is
\[
 Z_\delta:=D\Psi_\delta(U)[V]
 =\sum_i\lambda_i e^{ie\theta_i}
       \bigl(\zeta(v_i)+ie\beta_i\psi(v_i)\bigr).
\]
The first term has the scalar interpolation estimates just proved.
For the second, use \(\sum_i\lambda_i\beta_i=0\) to replace
\(e^{ie\theta_i}\psi(v_i)\) by
\(e^{ie\theta_i}\psi(v_i)-\psi(x)=O(h)\).
Since \(\beta_i,\partial_t\beta_i=O(h)\) and their spatial
gradients are bounded, the resulting sum is \(O(h^2)\), its time
derivative is \(O(h^2)\), and its spatial gradient is \(O(h)\).
These estimates are proportional to \(\|V\|_X\). In particular
\(Z_\delta-\zeta\) has those same three bounds. Its time derivative
includes all four terms
\[
 \partial_tZ_\delta=\sum_i\lambda_i e^{ie\theta_i}
 \left[\dot\zeta_i+ie\dot\beta_i\psi_i+ie\beta_i\dot\psi_i
       +ie\dot\theta_i(\zeta_i+ie\beta_i\psi_i)\right].
\]
No second time derivative is used.

The covariant quantities and their variations are
\begin{align*}
 Q_\delta&=\partial_t\Psi_\delta+ie\phi_\delta\Psi_\delta,&
 P_\delta&=\nabla\Psi_\delta-ieA_\delta\Psi_\delta,\\
 \delta Q_\delta&=\partial_tZ_\delta
       +ie(I_0\eta)\Psi_\delta+ie\phi_\delta Z_\delta,&
 \delta P_\delta&=\nabla Z_\delta
       -ie(I_1B)\Psi_\delta-ieA_\delta Z_\delta.
\end{align*}
They are uniformly bounded and differ from their continuum counterparts
by at most \(C_Rh\), with the variation bounds proportional to
\(\|V\|_X\). Finally the action density is
\(\tfrac12|E|^2-\tfrac12|B|^2+|Q|^2-|P|^2
-m^2|\psi|^2-g|\psi|^4/2\).
Its first variation is the sum of the corresponding bilinear pairings,
with matter potential contribution
\(-2(m^2+g|\psi|^2)\operatorname{Re}(\overline\psi\,\zeta)\).
On the bounded range just obtained both expressions are Lipschitz in their
arguments. Integrating their \(C_Rh\) differences over
\(I\times\Omega\), and using \(h\leq\delta\), proves the result.
The estimates hold almost everywhere for the stated Sobolev time regularity;
all spatial operations use the Lipschitz representatives on each macro-cell.
\end{proof}

Two controls delimit the argument. On the tetrahedron with vertices
\[
 0,\qquad(h,0,0),\qquad(0,h,0),\qquad(0,0,h),
\]
let \(A=\nabla(x_1^2)\).
The weighted original-potential path sum is \(x_1^2-hx_1\), whereas
the weighted Whitney-potential path sum is exactly zero. Interchanging
these two path definitions therefore changes the interpolant. For
a smooth bump \(b\) supported near an interior vertex with
\(b(0)=1\), the fields \(b_\epsilon(x)=b(x/\epsilon)\) satisfy
\(\|b_\epsilon\|_{H^1}^2=\epsilon^3\|b\|_{L^2}^2
+\epsilon\|\nabla b\|_{L^2}^2\to0\), while the nodal value stays
one. Thus this nodal construction cannot inherit an estimate for arbitrary
three-dimensional \(H^1\) data by point sampling.

\path{Lean/Screen/WhitneySpatialConsistency.lean} checks the finite
antisymmetric cancellation and error-decomposition algebra; the spatial
estimate itself is the analytic proof above. The executable refinement
probe in \path{code/electromagnetism/whitney_spatial_consistency.py}
uses manufactured nonzero-curvature fields on refinements of the same cone,
with independent quadrature and directional finite-difference checks.
These numerical checks test the implementation, not the uniform theorem
or a physical observer history. The controlled comparison retains a
supplied geometry, time coordinate and scalar-electrodynamics action.
